\subsection{Тестирование и анализ результатов}

\subsubsection{Тестирование корректности}

Тестирование корректности библиотеки осуществляется посредством
автоматических тестов, построенных на основе библиотеки Google
Test. Тесты организованы в три независимых набора, соответствующих
трём модулям библиотеки. Все тесты запускаются с включённым
Thread Sanitizer, что позволяет обнаруживать гонки данных
во время выполнения.

Набор тестов для \texttt{WorkStealingDeque} содержит 12 тестов,
разделённых на две группы. Первая группа проверяет корректность
операций в однопоточном режиме, вторая --- корректность при
конкурентном доступе нескольких потоков. Описание тестов
представлено в таблице~\ref{tab:tests-deque}.

\begin{table}[h]
\caption{Тесты корректности WorkStealingDeque}
\label{tab:tests-deque}
\begin{tabularx}{\textwidth}{|l|X|}
\hline
Тест & Что проверяется \\
\hline
\texttt{InitiallyEmpty} &
Очередь после создания пуста. \\
\hline
\texttt{PushIncreasesSize} &
Размер очереди увеличивается после \texttt{push}. \\
\hline
\texttt{PopOnEmptyReturnsFalse} &
\texttt{pop} возвращает \texttt{false} на пустой очереди. \\
\hline
\texttt{StealOnEmptyReturnsFalse} &
\texttt{steal} возвращает \texttt{false} на пустой очереди. \\
\hline
\texttt{PopIsLIFO} &
Владелец извлекает задачи в порядке LIFO. \\
\hline
\texttt{StealIsFIFO} &
Вор извлекает задачи в порядке FIFO. \\
\hline
\texttt{PopDecreasesSize} &
Размер очереди уменьшается после \texttt{pop}. \\
\hline
\texttt{GrowsWhenFull} &
Очередь корректно расширяется при переполнении. \\
\hline
\texttt{SingleThiefStealsAll} &
Один вор корректно крадёт все задачи. \\
\hline
\texttt{OwnerAndThiefConcurrently} &
Владелец и вор работают одновременно,
каждая задача выполняется ровно один раз. \\
\hline
\texttt{MultipleThievesConcurrently} &
Несколько воров работают одновременно,
каждая задача выполняется ровно один раз. \\
\hline
\texttt{SingleElementRaceCondition} &
Граничный случай: один элемент в очереди,
гонка между владельцем и вором. \\
\hline
\end{tabularx}
\end{table}

Тест \texttt{SingleElementRaceCondition} заслуживает отдельного
внимания. Он многократно (10\,000 итераций) создаёт очередь
с одним элементом и запускает одновременно владельца и вора,
каждый из которых пытается забрать этот элемент. После
завершения обоих потоков проверяется что задача была выполнена
ровно один раз. Данный тест выявляет ошибки в обработке
граничного случая, когда владелец и вор одновременно
претендуют на последний элемент очереди.

Набор тестов для \texttt{SimpleThreadPool} содержит 6 тестов,
проверяющих основные аспекты работы пула. Описание тестов
представлено в таблице~\ref{tab:tests-simple}.

\begin{table}[h]
\caption{Тесты корректности SimpleThreadPool}
\label{tab:tests-simple}
\begin{tabularx}{\textwidth}{|l|X|}
\hline
Тест & Что проверяется \\
\hline
\texttt{ReturnsCorrectResult} & Задача выполняется и возвращает корректный результат через \texttt{std::future}. \\
\hline
\texttt{MultipleTasksExecuted} & 100 задач выполняются и возвращают корректные результаты. \\
\hline
\texttt{AllTasksAreExecuted} & 1\,000 задач выполняются все без исключения. \\
\hline
\texttt{ExceptionPropagated} & Исключение внутри задачи корректно пробрасывается через \texttt{std::future::get()}. \\
\hline
\texttt{SingleThreadSequential} & При одном рабочем потоке задачи выполняются в порядке постановки. \\
\hline
\texttt{ThreadCountIsCorrect} & Метод \texttt{thread\_count()} возвращает корректное количество созданных потоков. \\
\hline
\end{tabularx}
\end{table}

Набор тестов для \texttt{WorkStealingPool} содержит 10 тестов.
По структуре он аналогичен набору для \texttt{SimpleThreadPool},
однако дополнительно включает тесты специфичные для механизма
work-stealing. Описание тестов представлено
в таблице~\ref{tab:tests-ws}.

\begin{table}[h]
\caption{Тесты корректности WorkStealingPool}
\label{tab:tests-ws}
\begin{tabularx}{\textwidth}{|l|X|}
\hline
Тест & Что проверяется \\
\hline
\texttt{ReturnsCorrectResult} &
Задача выполняется и возвращает корректный результат
через \texttt{std::future}. \\
\hline
\texttt{MultipleTasksExecuted} &
100 задач возвращают корректные результаты. \\
\hline
\texttt{AllTasksAreExecuted} &
1\,000 задач выполняются все без исключения. \\
\hline
\texttt{ExceptionPropagated} &
Исключение пробрасывается через \texttt{std::future}. \\
\hline
\texttt{ThreadCountIsCorrect} &
Метод \texttt{thread\_count()} возвращает корректное
число рабочих потоков. \\
\hline
\texttt{SingleThreadExecutes} &
Пул с одним потоком выполняет все задачи корректно. \\
\hline
\texttt{HeavyLoadAllTasks} &
100\,000 задач выполняются корректно под высокой
нагрузкой. \\
\hline
\texttt{UnevenTaskDurations} &
Все задачи выполняются при неравномерной нагрузке. \\
\hline
\texttt{ResultsUnderContention} &
Корректность результатов при конкуренции потоков. \\
\hline
\texttt{DestructorWaitsPending} &
Деструктор дожидается завершения всех задач. \\
\hline
\end{tabularx}
\end{table}

Результаты запуска тестов подтверждают корректность всех
реализованных компонентов: все 28 тестов из трёх наборов
успешно пройдены. Отсутствие сообщений об ошибках со стороны
Thread Sanitizer подтверждает корректность синхронизации
в реализованных структурах данных и отсутствие гонок данных
при конкурентном доступе нескольких потоков.